\documentclass[12pt]{article}
\usepackage{amsmath,amsfonts,fullpage,amsthm,graphicx}

\newtheorem*{theorem}{Theorem}
\newtheorem*{fact}{Fact}
\newtheorem*{goal}{Goal}
\newtheorem*{idea}{Idea}
\newtheorem*{defn}{Definition}
\newtheorem*{ex}{Example}

%Style section
%\setlength{\textheight}{10in}
%\setlength{\textwidth}{7in}
%\hoffset=-0.75in
 \pagestyle{plain}
% \setlength{\topmargin}{0in}
% \setlength{\headsep}{0in}
% \setlength{\oddsidemargin}{0in}
% \setlength{\evensidemargin}{0in}

\begin{document}


\pagestyle{empty} %use this to get rid of page numbers

\noindent
{Math 166 \hskip 1.5 truein  Names:\ \hrulefill}

\noindent
%\vskip 0.3truein
%\bigskip
%\noindent
{More shells}

%\noindent  \hskip 2.8 truein  Section Number:\ \hrulefill
\medskip

%\begin{goal} \rm
%Use integration to find the volume of various \emph{solids of revolution}.
%\end{goal}
%\smallskip
\noindent Graph the functions, draw a sample of a cross section of an arbitrary shell, and then compute.

\noindent(Hint: You will sometimes need to use multiple integrals.)

\begin{enumerate}

\item Use the shell method to find the volume of the solid obtained by revolving the region enclosed by $y = \sqrt{x+1}$, $y = 1$, $x = 0$, and $x=3$ about the $y$-axis.

\vspace{4in}
	
\item Use the shell method to find the volume of the solid obtained by revolving the regions enclosed by $y = x^2$, $y = 2-x$, $x = 0$, and $x=2$ about the $y$-axis.


\vspace{4in}

\item Use the shell method to find the volume of the solid obtained by revolving the region enclosed by $y = 1+\sqrt{x}$, $y = 1-\sqrt{x}$, $x = 0$, and $x=1$ about the $x$-axis.

\vspace{4in}

\item Use the shell method to find the volume of the solid obtained by revolving the regions enclosed by $y = x+1$, $y = 2 - x$, and $x = 0$  and $x=2$ about the line $y= -1$.
\newpage

\item Now use a combination of disks/washers and shells. Find the volume of the solid obtained by revolving the regions enclosed by $y = x+1$, $y = 2 - x$, and $x = 0$  and $x=2$ about the $y$-axis. A cross-section of this solid should be a wonky bow tie looking thing.

Use disks for the 'left' piece (inner solid) and shells for the 'right' piece (outer solid) or shells for the 'left', and washers for the 'right'.

\end{enumerate}

\end{document}
